\begin{theorem}[S-38: The residual tensor algebra and its reduced rank]
\label{thm:semireg-residual-tensor-rank}
Let \(X\) be a complex abelian variety and let \(T\subset\widehat X\) be a finite subgroup. Fix a coherent action by the tensor functors \[
 \Phi_t(-)=P_t\otimes(-),\qquad t\in T,
\] where \(P_t\) represents the indicated point of \(\widehat X\). A choice of the coherent multiplication maps makes \[
 \mathscr A=\bigoplus_{t\in T}P_t
 \tag{S38.1}
\] an associative strongly graded algebra; replacing every \(P_t\) by its inverse or passing to the opposite algebra only changes the conventions, not the conclusions. Let \(b:T\times T\to\mathbf C^*\) be its scalar commutator, let \(R=\operatorname{rad}(b)\), and put \[
 r=|R|,\qquad d=\sqrt{|T|/r}.
 \tag{S38.2}
\]

\textbf{(a) Center and reduced rank.} The integer \(d\) is positive. The center is \(\bigoplus_{t\in R}P_t\), its spectrum is a connected finite étale cover \[
 p:B\longrightarrow X,\qquad\deg p=r,
 \tag{S38.3}
\] and, after choosing an origin above zero, \(B\) is an abelian variety and \(p\) is an isogeny. Viewed over its center, \(\mathscr A\) is an Azumaya algebra \(\mathscr A_B\) of degree \(d\) on \(B\). A further finite étale isogeny trivializes \(\mathscr A_B\) as a matrix algebra.

Suppose a perfect object \(\mathcal E\) on \(X\) of positive virtual rank \(N\) has a coherent \(T\)-equivariant lift. Then \[
 m:=\frac{N}{rd}\in\mathbf Z_{>0},\qquad
 |T|r=(rd)^2\mid N^2.
 \tag{S38.4}
\] If \(\beta=[\mathscr A_B]\in\operatorname{Br}(B)\), then \[
 m\beta=0,\qquad d\beta=0,
 \qquad\operatorname{ind}(\beta_{\mathbf C(B)})\mid m.
 \tag{S38.5}
\] In particular, \(m=1\), or more generally \(\gcd(m,d)=1\), forces \(\beta=0\). Here the index is the usual index over the function field; no assertion that a minimal generic splitting rank is realized by a global bundle is needed.

\textbf{(b) Exact character conversion when the Brauer class splits.} If \(\beta=0\), choose a vector bundle \(S\) of rank \(d\) with \(\mathscr A_B\simeq\mathcal End(S)\). Then \[
 \operatorname{ch}(S)=d\exp\!\left(\frac{c_1(S)}d\right).
 \tag{S38.6}
\] There is an ordinary perfect object \(F\) of virtual rank \(m\) on \(B\) such that \[
 p^*\operatorname{ch}(\mathcal E)
       =rd\exp\!\left(\frac{c_1(S)}d\right)\operatorname{ch}(F).
 \tag{S38.7}
\] For the numerical conclusion it is enough to construct \(F\) from the alternating sum of the equivariant cohomology sheaves. Thus (S38.7) does not require an unproved comparison of equivariant derived categories.

\textbf{(c) Balanced real multiplication.} Suppose now that \(X\) is a polarized abelian fourfold with a Rosati-invariant rational real-quadratic action \(\mathbf Q(\sqrt{\Delta_0})\), of type \((2,2)\), where \(\Delta_0>1\) is squarefree. Use the notation of S-30, \[
 D=\int_X\frac{\theta^4}{24},\qquad
 C=\frac{\theta_1^2+\theta_2^2}{2},\qquad
 v_q=1-qC+q^2\frac{\theta^4}{24},\qquad q>0.
\] Assume \[
 e^{-c_1(\mathcal E)/N}\operatorname{ch}(\mathcal E)=Nv_q.
 \tag{S38.8}
\] If \(\beta=0\), the ordinary object \(F\) in (b), with the integral polarization \(\theta_B=p^*\theta\), satisfies \[
 e^{-c_1(F)/m}\operatorname{ch}(F)=m\,p^*v_q,
 \qquad D_B=rD.
 \tag{S38.9}
\]

If the original \(T\)-equivariant lift is simple, gluable and fully semiregular, the already checked S-22 numerical identity gives \[
 k=\frac{4DN^2q^2}{|T|}\in\{2,4\}.
 \tag{S38.10}
\] Then, whenever \(\beta=0\), \[
 \boxed{\Delta_0\mid m,\qquad k=2\Longrightarrow 2\mid m.}
 \tag{S38.11}
\] In particular \textbf{reduced rank \(m=1\) is impossible for every coherent pure-tensor action}, including nondegenerate commutator actions. This conclusion applies to positive virtual rank, not merely to vector bundles. It does not assume that the ordinary \(F\) is semiregular.
\end{theorem}
\begin{proof}
\par

\medskip\noindent\textbf{1. The center cover}\quad

The coherence maps identify \(P_t\otimes P_u\) with \(P_{t+u}\) and are associative. Comparing the maps for \((t,u)\) and \((u,t)\) gives a scalar \(b(t,u)\) because an automorphism of a line bundle on the connected projective variety \(X\) is scalar. Associativity shows that \(b\) is an alternating bicharacter. The standard finite Heisenberg representation theorem identifies every central-character block of the twisted group algebra with a matrix algebra of size \(d=\sqrt{|T|/r}\); see Polishchuk, Section 1.1, printed pp.4--5, with its reduction by the radical in Lemma 1.2.1, printed pp.6--7. These are imported background statements, not a new representation-theoretic claim.

The grading also identifies the center directly. A component in \(P_t\) commutes with every \(P_u\) exactly when \(b(t,u)=1\) for every \(u\). Since the graded components are independent, sums do not introduce further central terms. Hence \[
 Z(\mathscr A)=\bigoplus_{t\in R}P_t.
 \tag{S38.12}
\]

Let \(M\) be an exponent of \(T\). Pullback by \([M]:X\to X\) trivializes every \(P_t\). The resulting multiplication maps are constant, because they are holomorphic functions on a connected projective variety. Thus \([M]^*\mathscr A\) is a constant twisted group algebra and its center is the constant algebra \(\mathbf C^r\). It follows by finite étale descent that \(\operatorname{Spec}_X Z(\mathscr A)\) is a finite étale cover of degree \(r\), and that the corresponding algebra over this cover is étale locally \(\operatorname{Mat}_d\).

The cover is connected. For a nontrivial \(P\in\operatorname{Pic}^0(X)\) one has \(H^0(X,P)=0\): a nonzero section would give a numerically trivial effective divisor, whose intersection with a power of an ample class is positive unless the divisor is empty; the latter case trivializes \(P\). Therefore \[
 H^0(B,\mathcal O_B)=H^0(X,Z(\mathscr A))=\mathbf C.
\] Since \(B\) is proper and reduced, this proves connectedness. Analytically, a connected finite étale cover of \(X=V/\Lambda\) is \(V/\Lambda'\) for a finite-index subgroup \(\Lambda'\subset\Lambda\). Choosing an origin makes it a complex torus and the cover a group homomorphism; pullback of a polarization makes it an abelian variety. A connected component of the pullback of \(B\) along \([M]\) gives an isogeny \(f:Y\to B\) on which \(\mathscr A_B\) is a constant matrix algebra. This proves (a)'s geometric assertions.

\par

\medskip\noindent\textbf{2. Virtual rank and Brauer divisibility}\quad

Tensoring by a line bundle is exact for the usual coherent-sheaf \(t\)-structure. Thus every cohomology sheaf \(\mathcal H^i(\mathcal E)\) inherits its coherent tensor equivariance, equivalently a left \(\mathscr A\)-module structure. Passing to the center identifies it with a coherent \(\mathscr A_B\)-module \(M_i\) on \(B\), whose pushforward is \(\mathcal H^i(\mathcal E)\). Étale locally \(\mathscr A_B=\operatorname{Mat}_d\), and a module is a column module of rank \(d\) tensored with an ordinary coherent sheaf. Its generic rank is therefore \(d\) times an integer. Since \(B\) is connected, these ranks are well defined, and \[
 N=r\sum_i(-1)^i\operatorname{rk}(M_i)
   =rd\sum_i(-1)^i m_i.
\] This gives \(m=\sum_i(-1)^im_i=N/(rd)>0\) and (S38.4).

On the gerbe of local splittings of \(\mathscr A_B\), local Morita transforms of the \(M_i\) are twisted coherent sheaves of ranks \(m_i\), with the opposite Brauer sign if the alternative left/right convention is used. This gerbe is smooth over the regular variety \(B\), so these sheaves are perfect. Their alternating class has virtual rank \(m\). The determinant of that perfect class is an invertible \(m\)-fold twisted sheaf: a scalar \(z\) acts on it by \(z^m\), directly from the alternating determinant formula. Lieblich, Lemma 3.1.1.8 and Theorem A.3, therefore give \(m\beta=0\). The same determinant argument applied to the tautological local splitting bundle of rank \(d\) gives \(d\beta=0\); see also Stacks, Lemma 59.62.2.

At the generic point, the rank of every twisted coherent sheaf is divisible by the central division-algebra index, by the standard module classification, recorded in Lieblich, Proposition 3.1.2.1(ii). That index consequently divides the alternating rank \(m\). If \(\gcd(m,d)=1\), the two annihilation equations for \(\beta\) force \(\beta=0\). This proves all of (S38.5). Notice that no global vector bundle of rank equal to the generic index was asserted.

\par

\medskip\noindent\textbf{3. Splitting bundles have exponential character}\quad

Assume \(\beta=0\). The gerbe of splittings then has a global section, which is a vector bundle \(S\) of rank \(d\) with \(\mathcal End(S)\simeq\mathscr A_B\). Let \(f:Y\to B\) be the isogeny from Step 1. It gives \[
 \mathcal End(f^*S)\simeq\operatorname{Mat}_d(\mathcal O_Y)
                      =\mathcal End(\mathcal O_Y^{\oplus d}).
\] By Stacks, Lemma 59.62.1, there is a line bundle \(Q\) on \(Y\) such that \(f^*S\simeq Q^{\oplus d}\). Consequently \[
 f^*\operatorname{ch}(S)=d e^{c_1(Q)},\qquad
 f^*c_1(S)=d c_1(Q).
\] Pullback by a finite isogeny is injective in rational cohomology, since \(f_*f^*=(\deg f)\operatorname{id}\). This proves the full identity (S38.6), including its odd Chern-character components. It uses no classification of semihomogeneous bundles.

Ordinary Morita equivalence now writes \(M_i=S\otimes F_i\) for coherent sheaves \(F_i\) on \(B\). Since \(B\) is smooth, each \(F_i\) is perfect. Let \(F\) be the direct sum of the \(F_i\) placed in the cohomological degrees \(i\); its virtual rank is \(m\). Additivity of the Chern character and Riemann--Roch for the finite étale map \(p\), whose source and target have trivial tangent bundles, give \[
 \operatorname{ch}(\mathcal E)
    =p_*\bigl(\operatorname{ch}(S)\operatorname{ch}(F)\bigr).
\] For an isogeny, \(p^*p_*\) on rational cohomology is the sum of translations by the kernel, and each translation acts trivially; thus \(p^*p_*=r\operatorname{id}\). Inserting (S38.6) proves (S38.7). This cohomology-sheaf construction is sufficient even if no comparison functor for equivariant derived categories has been chosen.

\par

\medskip\noindent\textbf{4. Normalization and the arithmetic obstruction}\quad

Taking the degree-two component in (S38.7), and using \(N=rdm\), gives \[
 \frac{p^*c_1(\mathcal E)}N
      =\frac{c_1(S)}d+\frac{c_1(F)}m.
 \tag{S38.13}
\] Multiplying (S38.7) by the exponential of the negative of (S38.13), and applying (S38.8), proves (S38.9). The rational real-multiplication action transports through the isogeny, is Rosati-invariant for \(p^*\theta\), and has the same two eigenspace dimensions. The polarization \(p^*\theta\) is integral and its volume is \(D_B=rD\).

When the original equivariant numerical identity (S38.10) holds, its expression in the reduced variables is exactly \[
 k=\frac{4DN^2q^2}{|T|}
   =4D\frac{r^2d^2m^2q^2}{rd^2}
   =4D_Bm^2q^2.
 \tag{S38.14}
\] Apply the integral Chern-class conclusion S30.3 to the ordinary perfect object \(F\) in (S38.9). It yields \[
 m^4D_Bq^2=\frac{m^2k}{4}\in\Delta_0\mathbf Z_{>0}.
 \tag{S38.15}
\] If \(k=4\), squarefreeness gives \(\Delta_0\mid m\). If \(k=2\), integrality first forces \(m\) even, and the same divisibility gives \(\Delta_0\mid m\). These prove (S38.11). In particular \(m=1\) first forces \(\beta=0\) by Step 2 and then contradicts (S38.15), since \(\Delta_0>1\). No property of the full semiregularity map on \(F\) has entered the argument.
\end{proof}
\begin{remark}
The final numerical conclusion is for a pure-tensor action. The center and reduced-rank statements apply to the tensor subgroup of a larger action, but a nontrivial translation quotient requires a further equivariance argument. No Morita invariance of full semiregularity is assumed.\end{remark}
